\chapter{API design}

APIs are contracts.
Small design mistakes compound when clients and services depend on them.

\section{Resources and nouns}
Model resources as nouns and actions as methods.
Avoid verb-heavy endpoints like \code{/createTask}.

\section{Stable identifiers}
Use opaque, stable identifiers that do not leak internal state.
Do not overload identifiers with business meaning.

\section{Naming conventions}
Use consistent naming across resources and fields.
Prefer snake case in JSON fields if your ecosystem expects it, or camel case if that is standard.
Do not mix conventions within the same API.

\section{Nested resources}
Nested routes express ownership, but can become deep and repetitive.
Prefer one level of nesting and expose flat endpoints for cross-resource queries.

\section{Versioning}
Prefer backward-compatible changes and additive fields.
Use versioning when you must break clients.

\section{Versioning strategies}
\begin{tabularx}{\textwidth}{@{}lX@{}}
\toprule
Strategy & Tradeoff \\
\midrule
URI versioning & Easy to see, but splits clients across versions \\
Header versioning & Cleaner URLs, harder to debug \\
Query versioning & Simple but often abused \\
\bottomrule
\end{tabularx}

\section{Deprecation}
Deprecation is a process, not a flag.
Announce, measure usage, provide a migration path, and remove after a clear deadline.

\section{Error format}
Define one error shape across the API.

\begin{lstlisting}[language=JSON]
{
  "error": {
    "code": "validation_error",
    "message": "title is required",
    "field": "title"
  }
}
\end{lstlisting}

\section{Field naming}
Pick a naming convention and apply it everywhere.
If you use \code{created\_at} in one response, do not return \code{createdAt} in another.

\section{Filtering}
Filtering should be predictable and composable.
Use query parameters like \code{status=open} and \code{owner\_id=user\_7}.

\section{Expandable fields}
If a resource is large, allow clients to expand specific fields.
This avoids repeated follow-up requests while keeping defaults small.

\section{Partial responses}
For large collections, support selecting specific fields.
This reduces bandwidth and keeps responses stable under load.

\section{Idempotent creates}
If clients retry creates, accept an idempotency key.
Return the same resource if the create already succeeded.

\section{Pagination and filtering}
Use query parameters for filtering and sorting.
Keep pagination consistent across endpoints.

\section{Pagination example}
\begin{lstlisting}[language=JSON]
{
  "items": [
    {"id": "task_123", "title": "Write docs"}
  ],
  "next_cursor": "task_123"
}
\end{lstlisting}

\section{Filtering example}
\begin{lstlisting}[language=JSON]
{
  "filters": {
    "status": "open",
    "owner_id": "user_7"
  }
}
\end{lstlisting}

\section{Sorting}
Expose explicit sort fields and directions.
If sorting changes the meaning of pagination, use stable cursor-based pagination.

\section{Partial updates}
Define whether \code{PATCH} is merge semantics or JSON Patch.
Keep it consistent across endpoints.

\section{Long-running operations}
If an operation takes longer than a few seconds, return a job resource.
Clients should poll the job status instead of holding the request open.

\section{Job example}
\begin{lstlisting}[language=JSON]
{
  "job_id": "job_98",
  "status": "running",
  "result_url": "/jobs/job_98"
}
\end{lstlisting}

\section{Resource lifecycle}
Define lifecycle events for each resource.
For tasks: created, updated, completed, archived.
Ensure the API exposes timestamps for these transitions.

\section{Documentation}
Great APIs have short and accurate docs.
Document the request and response shape, error codes, and examples.

\section{Compatibility testing}
If multiple teams depend on your API, add contract tests.
They catch breaking changes before they reach production.

\section{Endpoint naming rules}
Keep endpoint names simple and stable:
\begin{itemize}[nosep]
  \item plural nouns for collections
  \item singular nouns for specific resources
  \item avoid verbs in paths
\end{itemize}

\section{Change policy}
Document how breaking changes are handled.
Common policies include a 90-day deprecation window and a migration guide.

\section{Search endpoints}
Search is different from filtering.
Expose a dedicated search endpoint when you need text search or relevance ranking.

\section{Bulk operations}
Bulk endpoints reduce round trips but increase complexity.
If you add bulk operations, make error handling explicit and predictable.

\section{Response envelopes}
If you wrap responses in an envelope, keep it consistent for all endpoints.
Do not mix raw and wrapped responses.

\section{Field validation}
Define validation rules once and reuse them across endpoints.
If you allow partial updates, validate only the provided fields but keep invariants.

\section{Pagination pitfalls}
Offset pagination is easy but unstable for changing datasets.
Cursor pagination is stable and scales better for large collections.

\section{Search response shape}
Search responses should include result counts and cursors when possible.
If counts are expensive, expose approximate counts explicitly.

\section{Actions and commands}
Some operations are not clean CRUD.
For actions like \code{/tasks/:id/complete}, prefer \code{POST} to an action subresource.
Keep the action names consistent and document the side effects.

\section{Write semantics}
Make the write semantics explicit so clients do not guess.

\begin{tabularx}{\textwidth}{@{}lX@{}}
\toprule
Method & Semantics \\
\midrule
PUT & Replace the full resource representation \\
PATCH & Apply a partial update with stable merge rules \\
POST & Create a new resource or trigger an action \\
\bottomrule
\end{tabularx}

\section{Default limits}
Pick default limits that protect the service:
\begin{itemize}[nosep]
  \item default page size, for example 50 items
  \item maximum page size, for example 200 items
  \item maximum payload size for create and update
\end{itemize}

If defaults are not stated, clients will assume the largest possible responses.

\section{Authorization scopes}
Define scopes or roles at the API level.
For example, a user may read tasks but not update them.
Document the scope requirements per endpoint and return consistent errors.

\section{Client SDKs and examples}
If the API is consumed by many teams, offer a small SDK or client library.
Even a thin wrapper helps standardize pagination, retries, and error handling.
Keep the SDK minimal and generated from the API contract when possible.

\section{Backward-compatible changes}
These changes are safe:
\begin{itemize}[nosep]
  \item adding a new field
  \item adding a new endpoint
  \item adding an optional query parameter
\end{itemize}

These changes are breaking:
\begin{itemize}[nosep]
  \item renaming fields
  \item changing field types
  \item changing pagination strategy
\end{itemize}

\section{OpenAPI as a contract}
An OpenAPI specification is a living contract.
Use it to generate documentation, validate requests, and drive tests.
Keep the spec in version control next to the code.

\section{Public versus internal APIs}
Public APIs need more stability and stricter compatibility guarantees.
Internal APIs can evolve faster, but they still need clear ownership.
Document which endpoints are public and which are internal only.

\section{Changelog and announcements}
APIs need a changelog.
Record additive changes, deprecations, and removals.
Even a short Markdown changelog reduces surprise for clients.

\section{Deprecation headers}
When an endpoint is deprecated, return a warning header and a link to the replacement.
This helps clients discover changes without reading every announcement.

\section{Error code taxonomy}
Define a small, stable set of error codes.
Avoid per-endpoint custom codes unless truly necessary.
Clients should be able to handle errors generically.

\section{Consistency across resources}
The same concept should look the same everywhere.
If \code{status} is an enum for tasks, do not expose it as a free-text field elsewhere.
Consistency reduces client complexity and mistakes.

\section{Response examples}
Include at least one example request and response per endpoint.
Examples are the fastest way for clients to integrate correctly.

\section{API review workflow}
Before a new API launches, run a short review:
\begin{itemize}[nosep]
  \item check naming and resource shape
  \item validate pagination and limits
  \item confirm error codes and edge cases
  \item verify backward compatibility
\end{itemize}

\section{Design workshop}
A short design workshop saves weeks of refactors.
Gather the backend and client teams and walk through:
\begin{itemize}[nosep]
  \item the primary user flows
  \item the minimal set of resources
  \item error cases and retries
  \item expected growth and versioning needs
\end{itemize}

Document the decisions and turn them into a contract.

\section{Contract testing loop}
Contract tests should run in CI and staging.
If a contract test fails, block the deploy until the change is reviewed.
This forces breaking changes to be explicit.

\section{API governance}
Governance does not mean bureaucracy.
It means a small set of rules that keep the API coherent:
\begin{itemize}[nosep]
  \item a single owner for each endpoint
  \item a consistent error schema
  \item a known deprecation process
\end{itemize}

\section{Client migration strategy}
When you must break compatibility, plan migration:
\begin{itemize}[nosep]
  \item announce the change and timeline
  \item provide a compatibility layer or adapter
  \item track adoption by client
  \item remove the old path on the announced date
\end{itemize}

\section{Field evolution playbook}
Fields evolve more often than endpoints.
Use a playbook to keep changes safe:
\begin{itemize}[nosep]
  \item add the new field and document it
  \item deploy clients that read the new field
  \item update writers to populate the new field
  \item deprecate and remove the old field
\end{itemize}

\section{API metrics}
Track API usage and error rates by endpoint and by client.
Metrics reveal which endpoints are critical and which can be changed safely.
Expose the top endpoints in a dashboard.

\section{Client types}
Different clients have different constraints.
Design with those constraints in mind:

\begin{tabularx}{\textwidth}{@{}lX@{}}
\toprule
Client & Constraint \\
\midrule
Mobile app & limited bandwidth, frequent retries \\
Browser & CORS and cookie behavior \\
Server-to-server & high throughput, strict SLAs \\
\bottomrule
\end{tabularx}

\section{Discovery and navigation}
Provide stable entry points for clients.
If you expose a root endpoint or API index, keep it backward compatible.
This reduces the risk of hard-coded URLs in clients.

\section{Case study: mobile client constraints}
Mobile clients retry more often due to unreliable networks.
They also cache aggressively to reduce bandwidth.
Design endpoints so retries are safe and responses are cacheable.
If the API is not retry-safe, mobile clients will create duplicate data.

\section{SDK versioning}
If you publish a client SDK, version it separately from the API.
The SDK should track breaking changes and provide deprecation warnings.
This helps clients update without surprises.

\section{Error handling examples}
Provide example error payloads for common failures:
\begin{itemize}[nosep]
  \item validation errors with field-level details
  \item authorization failures with a stable error code
  \item rate limit responses with retry hints
\end{itemize}

\section{Caching in API design}
Caching is a design decision.
If you expect heavy reads, include ETag and cache headers in the contract.
Document cache lifetimes so clients can rely on them.

\section{Case study: breaking change rollout}
Suppose you need to replace \code{status} with a richer object.
The safe approach is a staged rollout:
\begin{itemize}[nosep]
  \item add a new \code{status\_details} field
  \item update clients to read both fields
  \item backfill and populate the new field on write
  \item deprecate \code{status} after all clients migrate
\end{itemize}

This avoids hard breaks and gives clients a clear path forward.

\section{Compatibility testing example}
Contract tests can be simple:
\begin{itemize}[nosep]
  \item verify required fields in responses
  \item check status codes for error conditions
  \item ensure pagination fields are present
\end{itemize}

\begin{warningbox}{Warning: breaking changes by accident}
Renaming a field or changing its type is a breaking change.
If you are unsure, add a new field and deprecate the old one.
\end{warningbox}

\section{Example resource representation}
\begin{lstlisting}[language=JSON]
{
  "id": "task_123",
  "project_id": "proj_7",
  "title": "Write docs",
  "status": "open",
  "created_at": "2026-02-18T10:00:00Z",
  "updated_at": "2026-02-18T12:00:00Z"
}
\end{lstlisting}

\section{Review checklist}
Before releasing a new endpoint, review:
\begin{itemize}[nosep]
  \item status codes and error shapes
  \item pagination and limits
  \item authentication and authorization requirements
  \item backward compatibility with existing clients
\end{itemize}
\begin{patternbox}{Pattern: add fields, do not change meaning}
If you need new behavior, add a field instead of changing an existing one.
This keeps old clients working.
\end{patternbox}

\begin{checklistbox}{Checklist: API design}
\begin{itemize}[nosep]
  \item Use nouns for resources.
  \item Keep identifiers stable.
  \item Standardize error responses.
\end{itemize}
\end{checklistbox}
